\section{Arquitectura e implementación}
\label{sec:arquitectura_implementacion}

La implementación se desarrolló en Python siguiendo una arquitectura
modular que separa la especificación funcional de Keccak, la formulación
MILP, la búsqueda de trayectorias, la validación de testigos y la
generación de resultados.

Esta separación permite contrastar las soluciones del solver con una
implementación que no depende del modelo de optimización. De este modo,
el valor objetivo informado por CBC no se acepta de manera aislada, sino
que se verifica mediante la reconstrucción completa de los estados y de
sus diferencias.

PuLP se utilizó como biblioteca de modelado y CBC como solver de
programación lineal entera mixta
\cite{pulp_docs,cbc_guide}.


\subsection{Arquitectura general}
\label{subsec:arquitectura_general}

El flujo de trabajo se organiza en cinco etapas:

\begin{enumerate}
    \item implementar las transformaciones de Keccak reducido mediante
    funciones binarias de referencia;

    \item construir el modelo MILP exacto para dos ejecuciones simultáneas;

    \item incorporar las diferencias XOR, las variables de actividad y la
    función objetivo;

    \item resolver los experimentos mediante CBC o mediante una búsqueda
    restringida para tres rondas;

    \item reconstruir y validar los testigos con la implementación
    funcional independiente.
\end{enumerate}

El flujo puede resumirse como

\[
\text{Keccak de referencia}
\longrightarrow
\text{modelo MILP pareado}
\longrightarrow
\text{búsqueda de soluciones}
\longrightarrow
\text{reconstrucción}
\longrightarrow
\text{validación}.
\]

La arquitectura distingue claramente entre la obtención de una solución
y su comprobación. Esto reduce el riesgo de aceptar resultados producidos
por errores en la formulación, en la extracción de variables o en la
interpretación del estado del solver.


\subsection{Componentes principales}
\label{subsec:componentes_principales}

La Tabla~\ref{tab:componentes_implementacion} resume los componentes
principales del proyecto.

\begin{table}[htbp]
    \centering
    \caption{Componentes principales de la implementación.}
    \label{tab:componentes_implementacion}
    \begin{tabular}{p{0.31\textwidth}p{0.61\textwidth}}
        \hline
        \textbf{Componente} & \textbf{Responsabilidad} \\
        \hline

        Implementación de referencia
        &
        Ejecuta las transformaciones $\theta$, $\rho$, $\pi$, $\chi$ e
        $\iota$ directamente sobre arreglos binarios.
        \\[3pt]

        Modelo MILP de Keccak
        &
        Representa cada transformación mediante variables binarias,
        restricciones XOR, permutaciones y linealizaciones de productos.
        \\[3pt]

        Modelo diferencial pareado
        &
        Construye dos ejecuciones simultáneas y calcula las diferencias
        exactas en las fronteras de cada ronda.
        \\[3pt]

        Modelo de S-boxes activas
        &
        Incorpora las variables $a_{r,y,k}$, la función objetivo y las
        restricciones auxiliares para cotas, soportes y conteos por ronda.
        \\[3pt]

        Búsqueda restringida
        &
        Explora exhaustivamente trayectorias de tres rondas pertenecientes
        a la familia $2+2+c$.
        \\[3pt]

        Ejecutor experimental
        &
        Ejecuta los seis escenarios, almacena los resultados y genera los
        artefactos necesarios para su análisis.
        \\[3pt]

        Pruebas automatizadas
        &
        Verifican transformaciones, restricciones, testigos, búsquedas y
        archivos de salida.
        \\
        \hline
    \end{tabular}
\end{table}


\subsection{Implementación funcional de referencia}
\label{subsec:implementacion_referencia}

La implementación de referencia opera directamente sobre estados
binarios y reproduce el orden completo de una ronda:

\[
A_r
\xrightarrow{\theta}
T_r
\xrightarrow{\rho}
R_r
\xrightarrow{\pi}
B_r
\xrightarrow{\chi}
C_r
\xrightarrow{\iota}
A_{r+1}.
\]

Sus funciones se utilizan como una especificación independiente del
modelo MILP. En particular, permiten:

\begin{itemize}
    \item comprobar las transformaciones individuales;
    \item ejecutar una o varias rondas completas;
    \item reconstruir los estados obtenidos mediante el solver;
    \item recalcular las diferencias antes de cada aplicación de $\chi$;
    \item verificar el número de S-boxes activas por ronda.
\end{itemize}

La existencia de esta implementación independiente es importante porque
una solución MILP puede ser matemáticamente consistente con un modelo
incorrectamente programado. La comparación bit a bit permite detectar
este tipo de errores.


\subsection{Modelo MILP pareado}
\label{subsec:implementacion_modelo_pareado}

El modelo contiene dos copias de la permutación reducida:

\[
A_r^L
\qquad \text{y} \qquad
A_r^R.
\]

Cada ejecución satisface las restricciones exactas de las
transformaciones de Keccak. A partir de ambas copias se calculan las
diferencias

\[
\Delta A_r
=
A_r^L \oplus A_r^R,
\]

y

\[
\Delta B_r
=
B_r^L \oplus B_r^R.
\]

La clase que incorpora la minimización de actividad extiende este modelo
con las variables

\[
a_{r,y,k},
\]

que identifican las aplicaciones activas de $\chi$.

Además de la función objetivo global, la implementación permite:

\begin{itemize}
    \item imponer una cota sobre la actividad total;
    \item fijar el número de S-boxes activas en una ronda;
    \item fijar un soporte activo concreto;
    \item fijar los estados iniciales izquierdo y derecho;
    \item extraer los estados y soportes de una solución.
\end{itemize}

Estas operaciones se utilizan tanto para resolver los problemas de
optimización como para validar candidatos generados por otros
procedimientos.


\subsection{Búsqueda restringida para tres rondas}
\label{subsec:busqueda_restringida_implementacion}

La resolución directa del modelo global de tres rondas presenta un
crecimiento considerable del espacio combinatorio. Para obtener testigos
reproducibles se implementó una búsqueda restringida a trayectorias con
actividad

\[
2+2+c.
\]

La elección de esta familia parte del resultado exacto de dos rondas:
existen trayectorias globalmente óptimas con dos S-boxes activas en cada
ronda. La búsqueda analiza cómo pueden extenderse estas trayectorias hacia
una tercera ronda.

Para cada trayectoria compatible con actividad $2+2$, se enumeran las
realizaciones concretas de las dos S-boxes activas en la segunda ronda.
Cada S-box local contiene cinco bits y, por tanto, admite

\[
2^5=32
\]

valores absolutos posibles. Al existir dos S-boxes activas se evalúan

\[
32^2=1024
\]

realizaciones por trayectoria.

La búsqueda sigue el procedimiento:

\begin{enumerate}
    \item generar los soportes compatibles con actividad $2+2$;

    \item construir las realizaciones absolutas de las S-boxes activas;

    \item propagar cada realización hacia la tercera ronda;

    \item calcular el soporte de entrada a $\chi$ en dicha ronda;

    \item conservar los candidatos con menor actividad total;

    \item reconstruir los mejores candidatos mediante la implementación
    de referencia.
\end{enumerate}

El procedimiento es exhaustivo dentro de la familia $2+2+c$, pero no
dentro del espacio completo de trayectorias de tres rondas. Esta
distinción se mantiene en la interpretación de los resultados.


\subsection{Ejecución de los experimentos}
\label{subsec:ejecucion_experimentos}

El ejecutor principal procesa las configuraciones

\[
(z,R)
\in
\{
(4,1),
(4,2),
(4,3),
(8,1),
(8,2),
(8,3)
\}.
\]

La implementación dispone de dos modalidades:

\begin{description}
    \item[Modo rápido.] Utiliza los resultados permanentes de las
    búsquedas costosas y regenera las tablas, los testigos y las figuras.

    \item[Modo completo.] Repite la búsqueda restringida de tres rondas
    y reconstruye los resultados desde las trayectorias enumeradas.
\end{description}

Los resultados se almacenan en formatos estructurados para evitar la
transcripción manual de valores. Los artefactos principales contienen:

\begin{itemize}
    \item la matriz consolidada de los seis experimentos;
    \item los estados y soportes de los testigos;
    \item las cotas inferiores y superiores;
    \item la información del entorno de ejecución;
    \item la figura comparativa de resultados.
\end{itemize}

Los nombres y comandos concretos de reproducción se presentan en la
Sección~\ref{sec:reproducibilidad}.


\subsection{Estrategia de validación}
\label{subsec:estrategia_validacion}

La validación se realiza en cuatro niveles.


\subsubsection{Validación de transformaciones}

Las funciones de referencia se verifican mediante casos deterministas y
propiedades estructurales. Se comprueban las dimensiones de los estados,
la aritmética modular de los índices y la correspondencia entre las
transformaciones individuales y la ronda completa.


\subsubsection{Validación del modelo MILP}

Para estados fijados, se comparan las variables del modelo con la
implementación de referencia. La validación comprueba que:

\begin{enumerate}
    \item los estados intermedios coincidan bit a bit;
    \item las diferencias XOR sean correctas;
    \item la salida de $\chi$ corresponda a su definición booleana;
    \item las variables de actividad coincidan con los soportes de
    $\Delta B_r$;
    \item la suma de actividades coincida con el valor objetivo.
\end{enumerate}


\subsubsection{Validación de testigos}

Cada testigo contiene los estados iniciales

\[
A_0^L
\qquad \text{y} \qquad
A_0^R.
\]

A partir de estos estados se ejecutan nuevamente las rondas y se calculan
los soportes

\[
S_r
=
\left\{
(y,k):
\bigvee_{x=0}^{4}
\Delta B_r[x,y,k]=1
\right\}.
\]

El testigo se acepta únicamente cuando

\[
m_r=|S_r|
\]

coincide con la actividad almacenada para cada ronda y la suma

\[
\sum_{r=0}^{R-1}m_r
\]

coincide con la actividad total reportada.


\subsubsection{Validación automatizada}

La suite de pruebas cubre la implementación funcional, el modelo MILP,
las búsquedas, la reconstrucción de testigos y la generación de
artefactos.

En la versión utilizada para el informe se aprobaron

\[
245
\]

pruebas automatizadas. Este número se presenta como evidencia de cobertura
del proyecto, no como una demostración formal de corrección completa del
software.


\subsection{Criterio de aceptación de resultados}
\label{subsec:criterio_aceptacion_resultados}

Un resultado se incorpora a las conclusiones únicamente cuando satisface
las siguientes condiciones:

\begin{enumerate}
    \item el estado del solver permite una interpretación matemática
    válida;

    \item existe un testigo reconstruible cuando se afirma factibilidad;

    \item la implementación de referencia reproduce la actividad
    reportada;

    \item las cotas y los niveles de optimalidad se presentan sin exceder
    la evidencia disponible.
\end{enumerate}

Por tanto, una solución encontrada por CBC se interpreta como cota
superior mientras no exista una prueba de optimalidad. De forma análoga,
los mejores resultados de la búsqueda $2+2+c$ se consideran óptimos
restringidos y no mínimos globales de tres rondas.